\section{Panoramica Operativa e Strategie di Controllo}

L'architettura non convenzionale del velivolo, caratterizzata da una configurazione a tricottero con rotori anteriori avanzati e un rotore di coda a due gradi di libertà (2-DOF), richiede una gestione dinamica della stabilità differenziata per le diverse condizioni operative. La strategia di controllo è segmentata in tre fasi principali:

\subsection{Fase 1: Volo Verticale (Hovering)}
In questa configurazione, la portanza è generata interamente dai tre rotori orientati verticalmente ($\theta_{1,2,3} \approx 90^\circ$). La sfida principale risiede nel \textbf{bilanciamento del momento di imbardata}. Poiché i rotori anteriori sono controrotanti, la coppia di reazione residua prodotta dal motore posteriore deve essere compensata attivamente per evitare rotazioni indesiderate attorno all'asse $Z_b$.
\begin{itemize}
    \item \textbf{Strategia:} Si utilizza il grado di libertà di tilting laterale del rotore di coda ($\theta_3$) per generare una componente di spinta trasversale. Tale forza crea un momento di imbardata uguale e opposto alla coppia resistente del motore, garantendo l'equilibrio statico dell'assetto.
\end{itemize}

\subsection{Fase 2: Transizione}
Rappresenta la fase a maggiore criticità dinamica, in cui i rotori anteriori ruotano progressivamente verso la configurazione orizzontale per generare spinta propulsiva.
\begin{itemize}
    \item \textbf{Sfide:} Il sistema è soggetto a forti non linearità dovute all'interferenza aerodinamica dei flussi dei rotori sulle superfici alari e alla variazione del braccio di leva longitudinale del vettore spinta rispetto al baricentro.
    \item \textbf{Strategia:} Il controllo di beccheggio viene garantito dalla modulazione attiva della spinta del rotore posteriore, che compensa i momenti parassiti indotti dallo spostamento angolare dei motori anteriori.
\end{itemize}

\subsection{Fase 3: Volo Orizzontale (Crociera)}
Al raggiungimento della velocità di sostentamento ($v > v_{stall}$), l'ala assume il ruolo di superficie portante principale. Per massimizzare l'efficienza energetica e ridurre la resistenza aerodinamica, il sistema adotta una configurazione \textbf{sotto-attuata}.
\begin{itemize}
    \item \textbf{Strategia:} Il rotore posteriore viene disattivato ($\omega_3 = 0$). L'autorità di controllo è riallocata esclusivamente sui due rotori anteriori:
    \begin{itemize}
        \item \textbf{Rollio:} Gestito tramite tilt differenziale ($\theta_1 \neq \theta_2$).
        \item \textbf{Imbardata:} Gestita tramite modulazione differenziale della spinta ($\Delta \omega_{1,2}$).
        \item \textbf{Beccheggio:} Affidato alla stabilità intrinseca del profilo aerodinamico e a correzioni fini della spinta propulsiva.
    \end{itemize}
\end{itemize}